\begin{abstract} 
近年来，视觉-语言-动作（Vision-Language-Action, VLA）模型已成为机器人领域一种具有变革性的方案，它通过在端到端学习框架中融合视觉与语言输入，使机器人能够执行复杂任务。
尽管 VLA 模型展现出显著能力，但也引入了新的攻击面。本文系统评估了其鲁棒性。考虑到机器人执行任务的独特需求，
我们设计的攻击目标针对机器人系统内在的空间特性与功能特性。具体而言，我们提出了两个利用空间基础来扰乱机器人动作的非目标攻击目标，以及一个操纵机器人轨迹的目标攻击目标。此外，我们设计了一种对抗补丁生成方法，在相机视野中放置一个小型彩色补丁，从而能够在数字环境和物理环境中有效实施攻击。
实验结果表明，在一组模拟机器人任务中，任务成功率出现了显著下降，最高可降低 100\%，这凸显了当前 VLA 架构中存在的关键安全缺口。
通过揭示这些脆弱性并提出可操作的评估指标，我们推动了对基于 VLA 的机器人系统安全性的理解与提升，同时强调了在部署到物理世界之前持续发展鲁棒防御策略的必要性\footnote{\href{https://vlaattacker.github.io/}{项目主页}.}.
\end{abstract}
